La estructura de archivos del proyecto es sencilla y está organizada para facilitar el desarrollo y mantenimiento del código.
En la raíz del proyecto se encuentra el directorio `src/`, que contiene todo el código fuente.
Dentro de `src/`, hay un subdirectorio `main/` que alberga el código Java.

\vspace{0cm}
\begin{figure}[H]
    \centering
    \caption{Estructura de carpetas del proyecto}
    \vspace{-0.2cm}
    \fbox{
    \begin{minipage}{0.8\textwidth}
        \vspace{0.5cm}
        \dirtree{%
        .1 / (Raíz del proyecto).
        .2 src/.
        .3 main/.
        .4 java/.
        .5 Alumno.java \DTcomment{Clase Alumno}.
        .5 ClaseHorario.java \DTcomment{Clase Horario}.
        .5 Generador.java \DTcomment{Junta la lógica principal}.
        .5 HorarioUtils.java \DTcomment{funciones para manipular los horarios}.
        .5 Intervalo.java \DTcomment{Clase Intervalo}.
        .5 Main.java \DTcomment{Clase principal}.
        .5 PDFHorarioStripper.java \DTcomment{funciones para leer PDFs}.
        .5 SelectorDestino.java \DTcomment{Clase del selector de archivos}.
        .4 resources/pdf/ \DTcomment{Directorio para archivos PDF}.
        .5 *.pdf \DTcomment{Recursos PDF}.
        .2 target/ \DTcomment{Archivos compilados}.
        }
        \vspace{0.5cm}
    \end{minipage}
    }
    \label{fig:EstructuraCodigo}
\end{figure}

\subsubsection{Clase Alumno}
La clase `Alumno` funciona para almacenar datos y representa a un estudiante.
Contiene dos atributos: el nombre completo y el horario. Así como sus respectivos métodos para acceder a esta información.
Esta clase es fundamental para almacenar la información de cada alumno y sus preferencias de horario.
El código es:

\vspace{-0.2cm}
\begin{figure}[H]
    \centering
    \begin{minipage}{0.7\textwidth}
    \caption{Clase Alumno}
    \label{code:Alumno}
    \vspace{-0.2cm}
    \begin{lstlisting}[language=Java, basicstyle=\scriptsize\ttfamily]
public class Alumno {
    private final String nombreCompleto;
    private final Map<DayOfWeek, List<ClaseHorario>> horario;

    public Alumno(String nombreCompleto,
                  Map<DayOfWeek, List<ClaseHorario>> horario) {
        this.nombreCompleto = nombreCompleto;
        this.horario = horario;
    }

    public String getNombreCompleto() {
        return nombreCompleto;
    }

    public Map<DayOfWeek, List<ClaseHorario>> getHorario() {
        return horario;
    }
}
    \end{lstlisting}
    \end{minipage}
\end{figure}

Especificaciones:
\begin{itemize}
    \item Contiene el nombre completo del alumno.
    \item Contiene el horario. Está declarado como un Map donde la clave es la clase \texttt{DayOfWeek} que es el día de la semana
    y tiene el valor \texttt{List<ClaseHorario>} que almacena el horario para cada día de la semana. Por ejemplo:
    \begin{center}
        Monday $\rightarrow$ [Matemáticas 8:00, Física 10:00]
    \end{center}
\end{itemize}

\subsubsection{Clase Horario}
La clase `ClaseHorario` representa una materia específica dentro del horario de un alumno.
Contiene atributos para almacenar el código de la materia, el nombre, el día de la semana y
el intervalo de tiempo en el que se imparte la clase. Esta clase es esencial para organizar
y gestionar los horarios de los alumnos, permitiendo una fácil manipulación de la información
relacionada con las materias que cada alumno cursa.

\vspace{-0.2cm}
\begin{figure}[H]
    \centering
    \begin{minipage}{0.7\textwidth}
    \caption{Clase Horario}
    \label{code:Horario}
    \vspace{-0.2cm}
    \begin{lstlisting}[language=Java, basicstyle=\scriptsize\ttfamily]
public class ClaseHorario {
    private final String codigo;
    private String nombre;
    private final DayOfWeek dia;
    private final Intervalo intervalo;

    public ClaseHorario(String codigo, String nombre, DayOfWeek dia, Intervalo intervalo) {
        this.codigo = codigo;
        this.nombre = nombre;
        this.dia = dia;
        this.intervalo = intervalo;
    }

    public String getCodigo() {
        return codigo;
    }
    public String getNombreMateria() {
        return nombre;
    }
    public void setNombreMateria(String nombre) {
        this.nombre = nombre;
    }
    public DayOfWeek getDia() {
        return dia;
    }
    public Intervalo getIntervalo() {
        return intervalo;
    }
}
    \end{lstlisting}
    \end{minipage}
\end{figure}

Especificaciones:
\begin{itemize}
    \item Código: es la clave única de la materia. Por ejemplo: MAT101
    \item Nombre: nombre de la materia
    \item DayOfWeek: guarda el día de la materia. Por ejemplo: DayOfWeek.MONDAY3
    \item Intervalo: es una clase que almacena el intervalo donde la materia es tomada. Por ejemplo: 9:00-11:00
\end{itemize}

\subsubsection{Generador}
La clase `Generador` es el núcleo del sistema, encargada de coordinar la lectura de los archivos PDF,
procesar la información de los alumnos y generar el archivo de salida con los horarios.
Esta clase implementa la lógica principal del programa, incluyendo la gestión de hilos para permitir la
interrupción del proceso por parte del usuario, así como la organización y escritura de los horarios en un formato legible.
Además, maneja las excepciones que puedan surgir durante la lectura de los archivos PDF, asegurando que el programa se ejecute de manera robusta y eficiente.

\vspace{-0.2cm}
\begin{figure}[H]
    \centering
    \begin{minipage}{0.9\textwidth}
    \caption{Código de la Función Principal de Generador}
    \label{code:Generador}
    \vspace{-0.2cm}
    \begin{lstlisting}[language=Java, basicstyle=\scriptsize\ttfamily]
public class Generador {
    private static final Path CARPETA_PDF =
            Paths.get("src/main/resources/pdf/");
    private static final Path ARCHIVO_SALIDA =
            Paths.get("horarios.txt");
    public void ejecutar() throws Exception {
        System.out.println("El programa iniciara en 3 segundos.");
        System.out.println("Presiona [ENTER] en cualquier momento para cancelar...");

        // Hilo para vigilar el teclado y permitir la interrupcion del programa
        Thread hiloInterrupcion = new Thread(() -> {
            try {
                if (System.in.read() != -1) { 
                    System.out.println("\n[!] Ejecucion cancelada por el usuario.");
                    System.exit(0);
                }
            } catch (Exception e) {
            }
        });
        // Al ser Daemon, este hilo morira automaticamente cuando el programa termine
        hiloInterrupcion.setDaemon(true); 
        hiloInterrupcion.start();

        // El hilo principal hace la cuenta regresiva
        for (int i = 3; i > 0; i--) {
            System.out.print(i + "... ");
            Thread.sleep(1000);
        }
        System.out.println("\nIniciando procesamiento\n");

        Path rutaSalida = SelectorDestino.pedirRutaAlUsuario("horarios.txt");

        if (rutaSalida == null) {
        System.out.println("Proceso cancelado por el usuario.");
        return;
        }

        List<Alumno> alumnos;
        try {
            alumnos = cargarAlumnosDesdePDF();
        } catch (RuntimeException e) {
            System.err.println("[!] Proceso terminado.");
            return;
        }

        /* Se reemplazo ARCHIVO_SALIDA por rutaSalida */
        try (BufferedWriter writer = Files.newBufferedWriter(
                rutaSalida,
                StandardOpenOption.CREATE,
                StandardOpenOption.TRUNCATE_EXISTING
        )) {
            escribirHorariosAlumnos(alumnos, writer);
            writer.newLine();
            escribirHorariosLibres(alumnos, writer);
        }
    }
    ...
}
    \end{lstlisting}
    \end{minipage}
\end{figure}

La función \texttt{ejecutar()} (figura \ref{code:Generador}) es el punto de entrada del programa, donde se inicia la cuenta regresiva,
se configura el hilo de interrupción para permitir al usuario cancelar la ejecución, y se coordina la carga de los alumnos desde
los archivos PDF y la escritura de los horarios en el archivo de salida.
Esta función maneja las excepciones que puedan surgir durante el proceso,
asegurando que el programa se ejecute de manera robusta y eficiente.

\begin{figure}[H]
    \centering
    \begin{minipage}{0.9\textwidth}
    \caption{Código de la Función Cargar Alumnos de Generador}
    \label{code:Generador-CargarAlumnos}
    \vspace{-0.2cm}
    \begin{lstlisting}[language=Java, basicstyle=\scriptsize\ttfamily]
public class Generador {
    ...
    // =======================
    // Cargar alumnos
    // =======================
   private List<Alumno> cargarAlumnosDesdePDF() throws Exception {
        if (!Files.exists(CARPETA_PDF)) {
            throw new RuntimeException("No existe la carpeta de PDFs");
        }

        List<Alumno> alumnos = new ArrayList<>();

        try (DirectoryStream<Path> stream = Files.newDirectoryStream(CARPETA_PDF, "*.pdf")) {
            for (Path pdf : stream) {
                System.out.println("Procesando: " + pdf.getFileName());

                try {
                    try (PDDocument doc = PDDocument.load(pdf.toFile())) {
                        PDFHorarioStripper stripper = new PDFHorarioStripper();
                        stripper.getText(doc);
                        alumnos.add(new Alumno(stripper.getNombreAlumno(), stripper.getHorario()));
                    }
                } catch (IOException e) {
                    System.err.println("\n[!] Error fatal:");
                    System.err.println("[!] PDF no valido: " + pdf.getFileName());
                    System.err.println("[!] Detalle: " + e.getMessage());
                    System.err.println("[!] Se cancela toda la ejecucion.\n");
                    throw new RuntimeException(
                            "Proceso cancelado por archivo PDF no valido: " + pdf.getFileName(),
                            e
                    );
                }
            }
        }
        return alumnos;
    }
    ...
}
    \end{lstlisting}
    \end{minipage}
\end{figure}

La función \texttt{cargarAlumnosDesdePDF()} (figura \ref{code:Generador-CargarAlumnos}) genera
una lista de la clase \texttt{Alumno} que es responsable de leer los archivos PDF desde el directorio especificado,
envía cada archivo a la clase \texttt{PDFHorarioStripper} para extraer la información del alumno y su horario, y maneja las excepciones que puedan surgir durante este proceso.
Si se encuentra un archivo PDF no válido, se imprime un mensaje de error detallado y se cancela toda la ejecución del programa,
asegurando que el usuario esté informado sobre la causa del error.

\begin{figure}[H]
    \centering
    \begin{minipage}{0.9\textwidth}
    \caption{Código de la Función Escribir Horarios de Generador}
    \label{code:Generador-EscribirHorarios}
    \vspace{-0.2cm}
    \begin{lstlisting}[language=Java, basicstyle=\scriptsize\ttfamily]
public class Generador {
    ...
    // =======================
    // Escribir alumnos
    // =======================
    private void escribirHorariosAlumnos(
            List<Alumno> alumnos,
            BufferedWriter writer
    ) throws Exception {

        for (Alumno a : alumnos) {
            escribirAlumno(a, writer);
            writer.newLine();
        }
    }
    private void escribirAlumno(Alumno a, BufferedWriter writer)
            throws Exception {

        writer.write("[" + a.getNombreCompleto() + "]");
        writer.newLine();

        for (DayOfWeek dia : diasSemana()) {

            writer.write("  " + diaEnEspanol(dia) + ":");
            writer.newLine();

            var clases = a.getHorario().get(dia);

            if (clases == null || clases.isEmpty()) {
                writer.write("    N/H");
                writer.newLine();
                continue;
            }

            clases.sort(Comparator.comparing(
                    c -> c.getIntervalo().inicio
            ));

            for (ClaseHorario c : clases) {
                writer.write("    " +
                        c.getIntervalo().inicio + "-" +
                        c.getIntervalo().fin + "  " +
                        c.getNombreMateria()
                );
                writer.newLine();
            }
        }
    }
    ...
}
    \end{lstlisting}
    \end{minipage}
\end{figure}

La función \texttt{escribirHorariosAlumnos()} (figura \ref{code:Generador-EscribirHorarios}) es
responsable de escribir los horarios de cada alumno \textbf{en el archivo de salida}. Esta función itera sobre la lista de alumnos,
llamando a la función \texttt{escribirAlumno()} para cada uno, que se encarga de formatear y
escribir el nombre del alumno y su horario diario en un formato legible.

\begin{figure}[H]
    \centering
    \begin{minipage}{0.9\textwidth}
    \caption{Código de la Función Escribir Horarios Libres de Generador}
    \label{code:Generador-EscribirHorariosLibres}
    \vspace{-0.2cm}
    \begin{lstlisting}[language=Java, basicstyle=\scriptsize\ttfamily]
public class Generador {
    ...
    // =======================
    // Horarios libres
    // =======================
    private void escribirHorariosLibres(
            List<Alumno> alumnos,
            BufferedWriter writer
    ) throws Exception {

        writer.write("HORARIOS LIBRES COMUNES");
        writer.newLine();

        Map<DayOfWeek, List<Intervalo>> libres =
                HorarioUtils.horariosLibres(alumnos);

        for (DayOfWeek dia : diasSemana()) {

            writer.write("  " + diaEnEspanol(dia) + ":");
            writer.newLine();

            List<Intervalo> intervalos = libres.get(dia);

            if (intervalos == null || intervalos.isEmpty()) {
                writer.write("    N/H");
                writer.newLine();
                continue;
            }

            for (Intervalo i : intervalos) {
                writer.write("    " + i);
                writer.newLine();
            }
        }
    }
    ...
}
    \end{lstlisting}
    \end{minipage}
\end{figure}

La función \texttt{escribirHorariosLibres()} (figura \ref{code:Generador-EscribirHorariosLibres}) es responsable de escribir
los horarios libres comunes a todos los alumnos en el archivo de salida. Esta función utiliza la clase \texttt{HorarioUtils}
para calcular los intervalos de tiempo en los que ningún alumno tiene clases programadas, y luego formatea y
escribe esta información en el archivo de salida de manera legible para el usuario.

\begin{figure}[H]
    \centering
    \begin{minipage}{0.9\textwidth}
    \caption{Código de la Función Utilidades de Generador}
    \label{code:Generador-Utilidades}
    \vspace{-0.2cm}
    \begin{lstlisting}[language=Java, basicstyle=\scriptsize\ttfamily]
public class Generador {
    ...
    // =======================
    // Utilidades
    // =======================
    private List<DayOfWeek> diasSemana() {
        return List.of(
                DayOfWeek.MONDAY,
                DayOfWeek.TUESDAY,
                DayOfWeek.WEDNESDAY,
                DayOfWeek.THURSDAY,
                DayOfWeek.FRIDAY
        );
    }
    private String diaEnEspanol(DayOfWeek d) {
        return switch (d) {
            case MONDAY -> "Lunes";
            case TUESDAY -> "Martes";
            case WEDNESDAY -> "Miercoles";
            case THURSDAY -> "Jueves";
            case FRIDAY -> "Viernes";
            default -> "";
        };
    }
}
    \end{lstlisting}
    \end{minipage}
\end{figure}

Finalmente, la clase \texttt{Generador} también incluye funciones de utilidad para facilitar la manipulación de los datos.

\subsubsection{HorarioUtils}
La clase \texttt{HorarioUtils} es una colección de funciones estáticas que proporcionan utilidades para manipular y procesar los horarios de los alumnos.
Estas funciones incluyen el parseo de horarios desde texto, la unión de intervalos de tiempo, la identificación de horarios libres comunes entre los alumnos
y otras operaciones relacionadas con la gestión de los horarios.

\begin{figure}[H]
    \centering
    \begin{minipage}{0.9\textwidth}
    \caption{Código de la Función Parseo de Horarios de HorarioUtils}
    \label{code:HorarioUtils-Parseo}
    \vspace{-0.2cm}
    \begin{lstlisting}[language=Java, basicstyle=\scriptsize\ttfamily]
public class HorarioUtils {
    private HorarioUtils() {}
    // =======================
    // Parseo de horarios PDF
    // =======================
    public static Intervalo parseHorario(String texto) {

        String[] partes = texto.split("-");

        int hIni = Integer.parseInt(partes[0].substring(0, 2));
        int mIni = Integer.parseInt(partes[0].substring(2, 4));

        int hFin = Integer.parseInt(partes[1].substring(0, 2));
        int mFin = Integer.parseInt(partes[1].substring(2, 4));

        LocalTime inicio = LocalTime.of(hIni, mIni);

        LocalTime fin = (mFin == 59)
                ? LocalTime.of(hFin, 0).plusHours(1)
                : LocalTime.of(hFin, mFin);

        return new Intervalo(inicio, fin);
    }
    ...
}
    \end{lstlisting}
    \end{minipage}
\end{figure}

\begin{figure}[H]
    \centering
    \begin{minipage}{0.9\textwidth}
    \caption{Código de la Función Unión de Intervalos de HorarioUtils}
    \label{code:HorarioUtils-UnirIntervalos}
    \vspace{-0.2cm}
    \begin{lstlisting}[language=Java, basicstyle=\scriptsize\ttfamily]
public class HorarioUtils {
    ...
    // =======================
    // Unión de intervalos
    // =======================
    public static List<Intervalo>
    unirIntervalos(List<Intervalo> intervalos) {

        if (intervalos == null || intervalos.isEmpty())
            return List.of();

        intervalos.sort(
                Comparator.comparing(i -> i.inicio)
        );

        List<Intervalo> resultado = new ArrayList<>();
        Intervalo actual = intervalos.get(0);

        for (int i = 1; i < intervalos.size(); i++) {
            Intervalo sig = intervalos.get(i);

            if (!sig.inicio.isAfter(actual.fin)) {
                actual = new Intervalo(
                        actual.inicio,
                        actual.fin.isAfter(sig.fin)
                                ? actual.fin
                                : sig.fin
                );
            } else {
                resultado.add(actual);
                actual = sig;
            }
        }

        resultado.add(actual);
        return resultado;
    }
    ...
}
    \end{lstlisting}
    \end{minipage}
\end{figure}

\begin{figure}[H]
    \centering
    \begin{minipage}{0.9\textwidth}
    \caption{Código de la Función Unión de Horarios de HorarioUtils}
    \label{code:HorarioUtils-UnirHorarios}
    \vspace{-0.2cm}
    \begin{lstlisting}[language=Java, basicstyle=\scriptsize\ttfamily]
public class HorarioUtils {
    ...
    // =======================
    // Unión de horarios (alumnos)
    // =======================
    public static Map<DayOfWeek, List<Intervalo>>
    unirHorariosOcupados(List<Alumno> alumnos) {

        Map<DayOfWeek, List<Intervalo>> ocupados =
                new EnumMap<>(DayOfWeek.class);

        for (DayOfWeek d : DayOfWeek.values()) {
            ocupados.put(d, new ArrayList<>());
        }

        for (Alumno a : alumnos) {
            for (var entry : a.getHorario().entrySet()) {
                for (ClaseHorario c : entry.getValue()) {
                    ocupados.get(entry.getKey())
                            .add(c.getIntervalo());
                }
            }
        }

        Map<DayOfWeek, List<Intervalo>> resultado =
                new EnumMap<>(DayOfWeek.class);

        for (DayOfWeek dia : List.of(
                DayOfWeek.MONDAY,
                DayOfWeek.TUESDAY,
                DayOfWeek.WEDNESDAY,
                DayOfWeek.THURSDAY,
                DayOfWeek.FRIDAY
        )) {
            resultado.put(
                    dia,
                    unirIntervalos(ocupados.get(dia))
            );
        }

        return resultado;
    }
    ...
}
    \end{lstlisting}
    \end{minipage}
\end{figure}

\begin{figure}[H]
    \centering
    \begin{minipage}{0.9\textwidth}
    \caption{Código de la Función Complemento de Intervalos de HorarioUtils}
    \label{code:HorarioUtils-ComplementoIntervalos}
    \vspace{-0.2cm}
    \begin{lstlisting}[language=Java, basicstyle=\scriptsize\ttfamily]
public class HorarioUtils {
    ...
    // =======================
    // Complemento de intervalos
    // =======================
    public static List<Intervalo>
    complemento(List<Intervalo> ocupados,
                LocalTime inicio,
                LocalTime fin) {

        List<Intervalo> libres = new ArrayList<>();
        LocalTime cursor = inicio;

        for (Intervalo o : ocupados) {

            if (cursor.isBefore(o.inicio)) {
                libres.add(new Intervalo(cursor, o.inicio));
            }

            if (o.fin.isAfter(cursor)) {
                cursor = o.fin;
            }
        }

        if (cursor.isBefore(fin)) {
            libres.add(new Intervalo(cursor, fin));
        }

        return libres;
    }
    ...
}
    \end{lstlisting}
    \end{minipage}
\end{figure}

\begin{figure}[H]
    \centering
    \begin{minipage}{0.9\textwidth}
    \caption{Código de la Función Horarios Libres de HorarioUtils}
    \label{code:HorarioUtils-HorariosLibresComunes}
    \vspace{-0.2cm}
    \begin{lstlisting}[language=Java, basicstyle=\scriptsize\ttfamily]
public class HorarioUtils {
    ...
    // =======================
    // Horarios libres comunes
    // =======================
    public static Map<DayOfWeek, List<Intervalo>>
    horariosLibres(List<Alumno> alumnos) {

        Map<DayOfWeek, List<Intervalo>> ocupados =
                unirHorariosOcupados(alumnos);

        Map<DayOfWeek, List<Intervalo>> libres =
                new EnumMap<>(DayOfWeek.class);

        LocalTime INICIO = LocalTime.of(11, 0);
        LocalTime FIN = LocalTime.of(19, 0);

        for (DayOfWeek dia : ocupados.keySet()) {

            List<Intervalo> libresDia =
                    complemento(
                            ocupados.get(dia),
                            INICIO,
                            FIN
                    );

            libres.put(dia, libresDia);
        }

        return libres;
    }
}
    \end{lstlisting}
    \end{minipage}
\end{figure}

\subsubsection{Clase Intervalo}
La clase \texttt{Intervalo} es una clase simple que representa un intervalo de tiempo con una hora de inicio y una hora de fin.
Esta clase es utilizada para almacenar los intervalos de tiempo en los que se imparten las clases, así como para calcular los horarios libres comunes entre los alumnos.
Contiene dos atributos: \texttt{inicio} y \texttt{fin}, ambos de tipo \texttt{LocalTime} y un método \texttt{toString()} para representar el intervalo en formato legible.

\begin{figure}[H]
    \centering
    \begin{minipage}{0.9\textwidth}
    \caption{Código de la clase Intervalo}
    \label{code:Intervalo}
    \vspace{-0.2cm}
    \begin{lstlisting}[language=Java, basicstyle=\scriptsize\ttfamily]
import java.time.LocalTime;

public class Intervalo {

    public final LocalTime inicio;
    public final LocalTime fin;

    public Intervalo(LocalTime inicio, LocalTime fin) {
        this.inicio = inicio;
        this.fin = fin;
    }

    @Override
    public String toString() {
        return inicio + - + fin;
    }
}
    \end{lstlisting}
    \end{minipage}
\end{figure}

Especificaciones:
\begin{itemize}
    \item Contiene una hora final y una hora inicial, esto para intervalo en clase horario.
\end{itemize}

\subsubsection{Clase Main}
La clase \texttt{Main} es la clase principal del programa, que contiene el método \texttt{main()} que es el punto de entrada de la aplicación.
En este método se crea una instancia de la clase \texttt{Generador} y se llama a su método \texttt{ejecutar()} para iniciar el proceso de generación de los horarios.

\begin{figure}[H]
    \centering
    \begin{minipage}{0.9\textwidth}
    \caption{Código de la clase Main}
    \label{code:Main}
    \vspace{-0.2cm}
    \begin{lstlisting}[language=Java, basicstyle=\scriptsize\ttfamily]
public class Main {
    public static void main(String[] args) throws Exception {
        new Generador().ejecutar();
    }
}
    \end{lstlisting}
    \end{minipage}
\end{figure}

\subsubsection{Clase PDFHorarioStripper}
La clase \texttt{PDFHorarioStripper} es una extensión de la clase \texttt{PDFTextStripper} de la biblioteca Apache PDFBox,
diseñada para extraer información específica de los archivos PDF que contienen los horarios de los alumnos.
Esta clase implementa la lógica para identificar y extraer el nombre del alumno, los códigos de las materias, los nombres de las materias,
los días de la semana y los horarios de cada clase.

\begin{figure}[H]
    \centering
    \begin{minipage}{0.9\textwidth}
    \caption{Código de las constantes de la clase PDFHorarioStripper}
    \label{code:PDFHorarioStripper-Constantes}
    \vspace{-0.2cm}
    \begin{lstlisting}[language=Java, basicstyle=\scriptsize\ttfamily]
public class PDFHorarioStripper extends PDFTextStripper {
    // Constantes
    // =======================
    private static final Map<String, DayOfWeek> DIAS = Map.of(
            "LUNES", DayOfWeek.MONDAY,
            "MARTES", DayOfWeek.TUESDAY,
            "MIERCOLES", DayOfWeek.WEDNESDAY,
            "MIÉRCOLES", DayOfWeek.WEDNESDAY,
            "JUEVES", DayOfWeek.THURSDAY,
            "VIERNES", DayOfWeek.FRIDAY,
            "SABADO", DayOfWeek.SATURDAY,
            "SÁBADO", DayOfWeek.SATURDAY,
            "DOMINGO", DayOfWeek.SUNDAY
    );
    private static final String REGEX_CODIGO = "[A-Z]{3,4}-\\d{3}";
    private static final String REGEX_CODIGO_PARCIAL = "[A-Z]{3,4}-";
    private static final String REGEX_HORARIO_COMPLETO = "\\d{4}-\\d{4}";
    private static final String REGEX_HORARIO_INICIO = "\\d{4}-";
    private static final String REGEX_HORARIO_FIN = "\\d{4}";
    ...
}
    \end{lstlisting}
    \end{minipage}
\end{figure}

En la figura \ref{code:PDFHorarioStripper-Constantes} se pueden observar las constantes utilizadas en la clase \texttt{PDFHorarioStripper},
incluyendo un mapa para convertir los nombres de los días de la semana en español a objetos \texttt{DayOfWeek} y varias expresiones regulares
para identificar los códigos de las materias y los horarios en el texto extraído de los archivos PDF.

\begin{figure}[H]
    \centering
    \begin{minipage}{0.9\textwidth}
    \caption{Código de estado de la clase PDFHorarioStripper}
    \label{code:PDFHorarioStripper-Estado}
    \vspace{-0.2cm}
    \begin{lstlisting}[language=Java, basicstyle=\scriptsize\ttfamily]
public class PDFHorarioStripper extends PDFTextStripper {
    ...
    // =======================
    // Estado
    // =======================

    private final TreeMap<Float, DayOfWeek> columnasDias = new TreeMap<>();
    private final Map<DayOfWeek, List<ClaseHorario>> horario =
            new EnumMap<>(DayOfWeek.class);

    private String nombreAlumno;

    private String codigoMateriaActual;
    private String codigoParcial;

    private StringBuilder nombreMateriaActual;
    private boolean aceptandoNombreMateria;

    private String horarioParcial;

    private ClaseHorario ultimaClaseCreada;

    public PDFHorarioStripper() throws IOException {
        setSortByPosition(true);
        for (DayOfWeek d : DayOfWeek.values()) {
            horario.put(d, new ArrayList<>());
        }
    }

    ...
}
    \end{lstlisting}
    \end{minipage}
\end{figure}

En la figura \ref{code:PDFHorarioStripper-Estado} se pueden observar los atributos de estado de la clase \texttt{PDFHorarioStripper},
que se utilizan para almacenar la información extraída de los archivos PDF, así como para mantener el estado durante el proceso de extracción.
Estos atributos incluyen un mapa para registrar las columnas de los días de la semana, un mapa para almacenar el horario del alumno,
variables para el nombre del alumno, el código de la materia actual, el nombre de la materia actual y otras variables auxiliares para el proceso de extracción.

\begin{figure}[H]
    \centering
    \begin{minipage}{0.9\textwidth}
    \caption{Código de la clase PDFHorarioStripper}
    \label{code:PDFHorarioStripper}
    \vspace{-0.2cm}
    \begin{lstlisting}[language=Java, basicstyle=\scriptsize\ttfamily]
public class PDFHorarioStripper extends PDFTextStripper {
    ...
    @Override
    protected void writeString(String text, List<TextPosition> positions) {
        if (text == null || positions.isEmpty()) return;

        String limpio = text
                .replace('\u00A0', ' ')
                .trim()
                .toUpperCase();

        if (limpio.isEmpty()) return;

        TextPosition p = positions.get(0);
        float x = p.getXDirAdj();
        // =======================
        // Nombre del alumno
        // =======================
        if (limpio.startsWith("NOMBRE:")) {
            nombreAlumno = limpio.substring("NOMBRE:".length()).trim();
            return;
        }
        // =======================
        // Encabezados de días
        // =======================
        if (DIAS.containsKey(limpio)) {
            registrarColumnaDia(x, DIAS.get(limpio));
            return;
        }
        // =======================
        // Código de materia completo
        // =======================
        if (limpio.matches(REGEX_CODIGO)) {
            codigoMateriaActual = limpio;
            codigoParcial = null;
            nombreMateriaActual = new StringBuilder();
            ultimaClaseCreada = null;
            aceptandoNombreMateria = true;
            return;
        }
        // =======================
        // Código de materia partido
        // =======================
        if (limpio.matches(REGEX_CODIGO_PARCIAL)) {
            codigoParcial = limpio;
            return;
        }
        if (codigoParcial != null && limpio.matches("\\d{3}")) {
            codigoMateriaActual = codigoParcial + limpio;
            codigoParcial = null;
            nombreMateriaActual = new StringBuilder();
            ultimaClaseCreada = null;
            aceptandoNombreMateria = true;
            return;
        }
        ...
    }

    ...
}
    \end{lstlisting}
    \end{minipage}
\end{figure}

En la figura \ref{code:PDFHorarioStripper} se puede observar el método \texttt{writeString()} de la clase \texttt{PDFHorarioStripper},
que es el método principal encargado de procesar cada fragmento de texto extraído del archivo PDF. Este método implementa la lógica para identificar y extraer el nombre del alumno,
los códigos de las materias, los nombres de las materias, los días de la semana y los horarios de cada clase, utilizando las constantes
y el estado definidos en la clase para guiar el proceso de extracción.

\begin{figure}[H]
    \centering
    \begin{minipage}{0.9\textwidth}
    \caption{Código de la clase PDFHorarioStripper}
    \label{code:PDFHorarioStripper-Horarios}
    \vspace{-0.2cm}
    \begin{lstlisting}[language=Java, basicstyle=\scriptsize\ttfamily]
public class PDFHorarioStripper extends PDFTextStripper {
    @Override
    protected void writeString(String text, List<TextPosition> positions) {
        ...
        // =======================
        // Nombre de materia (solo ANTES del horario)
        // =======================
        if (aceptandoNombreMateria
                && Character.isLetter(limpio.charAt(0))
                && !DIAS.containsKey(limpio)
                && !limpio.startsWith("TOTAL")
                && !limpio.startsWith("NOTAS")
                && !limpio.startsWith("VERSIÓN")
                && !limpio.matches(REGEX_CODIGO)
                && !limpio.matches(REGEX_CODIGO_PARCIAL)
                && !limpio.matches(REGEX_HORARIO_COMPLETO)
                && !limpio.matches(REGEX_HORARIO_INICIO)
                && !limpio.matches(REGEX_HORARIO_FIN)) {

            if (nombreMateriaActual.length() > 0)
                nombreMateriaActual.append(" ");

            String limpioMateria = limpiarNombreMateria(limpio);

            if (!limpioMateria.isEmpty()) {
                if (nombreMateriaActual.length() > 0)
                    nombreMateriaActual.append(" ");

                nombreMateriaActual.append(limpioMateria);
            }

            return;
        }
        // =======================
        // Horario completo
        // =======================
        if (limpio.matches(REGEX_HORARIO_COMPLETO)) {
            procesarHorario(limpio, x);
            return;
        }
        // =======================
        // Horario partido (inicio)
        // =======================
        if (limpio.matches(REGEX_HORARIO_INICIO)) {
            horarioParcial = limpio;
            return;
        }
        // =======================
        // Horario partido (fin)
        // =======================
        if (horarioParcial != null && limpio.matches(REGEX_HORARIO_FIN)) {
            String completo = horarioParcial + limpio;
            horarioParcial = null;
            procesarHorario(completo, x);
        }
    }

    ...
}
    \end{lstlisting}
    \end{minipage}
\end{figure}

En la figura \ref{code:PDFHorarioStripper-Horarios} se puede observar la parte del método \texttt{writeString()} que se encarga de procesar los horarios de las clases.
Esta sección del código identifica los fragmentos de texto que corresponden a los horarios, ya sea en formato completo o partido y llama a la función
\texttt{procesarHorario()} para extraer la información del horario y almacenarla en el estado de la clase.

\begin{figure}[H]
    \centering
    \begin{minipage}{0.9\textwidth}
    \caption{Código de la función procesarHorario de PDFHorarioStripper}
    \label{code:PDFHorarioStripper-ProcesarHorario}
    \vspace{-0.2cm}
    \begin{lstlisting}[language=Java, basicstyle=\scriptsize\ttfamily]
public class PDFHorarioStripper extends PDFTextStripper {
    ...
    // =======================
    // Procesamiento de horario
    // =======================

    private void procesarHorario(String texto, float x) {

        if (codigoMateriaActual == null) return;

        DayOfWeek dia = diaPorColumna(x);
        if (dia == null) return;

        Intervalo intervalo = HorarioUtils.parseHorario(texto);

        String nombre =
                (nombreMateriaActual != null && nombreMateriaActual.length() > 0)
                        ? normalizarNombreMateria(nombreMateriaActual.toString())
                        : null;

        ClaseHorario clase = new ClaseHorario(
                codigoMateriaActual,
                nombre,
                dia,
                intervalo
        );

        horario.get(dia).add(clase);
        ultimaClaseCreada = clase;

        // Después del primer horario ya NO se acepta más nombre
        aceptandoNombreMateria = false;
    }
    ...
}
    \end{lstlisting}
    \end{minipage}
\end{figure}

En la figura \ref{code:PDFHorarioStripper-ProcesarHorario} se puede observar la función \texttt{procesarHorario()} de la clase \texttt{PDFHorarioStripper},
que es responsable de procesar un fragmento de texto que representa un horario, extraer la información relevante y almacenarla en el estado de la clase.
Esta función verifica que haya un código de materia actual, identifica el día de la semana correspondiente a la posición del texto, parsea el horario utilizando la clase \texttt{HorarioUtils}
y luego crea una instancia de \texttt{ClaseHorario} que se agrega al horario del alumno.

\begin{figure}[H]
    \centering
    \begin{minipage}{0.9\textwidth}
    \caption{Código de las Utilidades - PDFHorarioStripper}
    \label{code:PDFHorarioStripper-Utilidades}
    \vspace{-0.2cm}
    \begin{lstlisting}[language=Java, basicstyle=\scriptsize\ttfamily]
public class PDFHorarioStripper extends PDFTextStripper {
    ...
    // =======================
    // Utilidades
    // =======================

    private void registrarColumnaDia(float x, DayOfWeek dia) {
        final float TOLERANCIA = 5f;
        for (float px : columnasDias.keySet()) {
            if (Math.abs(px - x) < TOLERANCIA) return;
        }
        columnasDias.put(x, dia);
    }

    private DayOfWeek diaPorColumna(float xHorario) {
        float mejor = Float.MAX_VALUE;
        DayOfWeek dia = null;

        for (var e : columnasDias.entrySet()) {
            float dx = Math.abs(e.getKey() - xHorario);
            if (dx < mejor) {
                mejor = dx;
                dia = e.getValue();
            }
        }
        return dia;
    }
    private String limpiarNombreMateria(String texto) {

        // Normaliza espacios
        texto = texto.replaceAll("\\s+", " ").trim();

        // Si no hay guion, no hay docente
        if (!texto.contains("-")) {
            return texto;
        }

        String[] partes = texto.split("-");

        if (partes.length < 2) {
            return texto;
        }

        String izquierda = partes[0].trim();   // Materia + Docente
        String derecha = partes[1].trim();     // Grupo

        // Quitar el último "token largo" de la izquierda (docente)
        String[] palabras = izquierda.split(" ");

        StringBuilder materia = new StringBuilder();
        for (String p : palabras) {
            // Heurística: apellidos suelen ser largos
            if (p.length() >= 12) {
                break;
            }
            materia.append(p).append(" ");
        }

        // Añadir grupo (ETIC, A, B, etc.)
        materia.append(derecha);

        return materia.toString().trim();
    }
    private String normalizarNombreMateria(String texto) {

        // Normaliza espacios
        texto = texto.replaceAll("\\s+", " ").trim();

        // Caso: MATERIA  DOCENTE - GRUPO
        // Ej: SEM DE TITULACION Y NETZAHUALCOYOTZI - ETIC
        return texto.replaceAll(
                "^(.+?)\\s+[A-ZÁÉÍÓÚÑ]{6,}(?:\\s+[A-ZÁÉÍÓÚÑ]{6,})*\\s*-\\s*([A-Z0-9]{1,6})$",
                "$1 $2"
        ).trim();
    }
    ...
}
    \end{lstlisting}
    \end{minipage}
\end{figure}

En la figura \ref{code:PDFHorarioStripper-Utilidades} se pueden observar varias funciones de utilidad dentro de la clase \texttt{PDFHorarioStripper},
incluyendo funciones para registrar las columnas de los días de la semana, identificar el día de la semana correspondiente a una posición dada, limpiar y
normalizar los nombres de las materias.

\begin{figure}[H]
    \centering
    \begin{minipage}{0.9\textwidth}
    \caption{Código de los getters - PDFHorarioStripper}
    \label{code:PDFHorarioStripper-Getters}
    \vspace{-0.2cm}
    \begin{lstlisting}[language=Java, basicstyle=\scriptsize\ttfamily]
    // =======================
    // Getters
    // =======================
    public Map<DayOfWeek, List<ClaseHorario>> getHorario() {
        return horario;
    }
    public String getNombreAlumno() {
        return nombreAlumno;
    }
    \end{lstlisting}
    \end{minipage}
\end{figure}

Finalmente, en la figura \ref{code:PDFHorarioStripper-Getters} se pueden observar los getters de la clase \texttt{PDFHorarioStripper}
que permiten acceder al nombre del alumno y su horario después de que se ha procesado el archivo PDF.

\subsubsection{Clase SelectorDestino}
La clase \texttt{SelectorDestino} es unicamente responsable de mostrar un diálogo al usuario para seleccionar la ubicación y
el nombre del archivo de salida donde se guardarán los horarios generados.

\begin{figure}[H]
    \centering
    \begin{minipage}{0.9\textwidth}
    \caption{Código de la clase SelectorDestino}
    \label{code:SelectorDestino}
    \vspace{-0.2cm}
    \begin{lstlisting}[language=Java, basicstyle=\scriptsize\ttfamily]
public class SelectorDestino {
    public static Path pedirRutaAlUsuario(String nombrePorDefecto) {
        // Intentar que el diálogo use la estética del SO actual (Windows, Linux, etc.)
        try {
            UIManager.setLookAndFeel(UIManager.getSystemLookAndFeelClassName());
        } catch (Exception e) {
            // Si falla, Java usará su tema visual por defecto
        }

        JFileChooser selector = new JFileChooser();
        selector.setDialogTitle("Seleccionar destino de guardado");
        selector.setSelectedFile(new File(nombrePorDefecto));
        // Filtro para asegurar que sea un archivo de texto
        FileNameExtensionFilter filtro = new FileNameExtensionFilter("Documento de texto (*.txt)", "txt");
        selector.setFileFilter(filtro);

        int resultado = selector.showSaveDialog(null);

        if (resultado == JFileChooser.APPROVE_OPTION) {
            File archivo = selector.getSelectedFile();
            // Verificación manual de extensión .txt si el usuario no la escribió
            String ruta = archivo.getAbsolutePath();
            if (!ruta.toLowerCase().endsWith(".txt")) {
                archivo = new File(ruta + ".txt");
            }
            return archivo.toPath();
        }
        return null; // El usuario cerró la ventana o canceló
    }
}
    \end{lstlisting}
    \end{minipage}
\end{figure}